\chapter{Choosing a Method: Decision Framework and Testbed Evidence}
\label{ch:selection}

Chapters~\ref{ch:success} through~\ref{ch:censored} each argued for a method inside its own regime of
observability. This chapter is the exit interview. Given a dataset, a question, and a tolerance for
assumptions: which pipeline do we run, what does it estimate, how do we attach uncertainty, and what is most
likely to go wrong? We compress the manual into two decision tables (Section~\ref{sec:sel-decisions}), then
present the measured evidence from the synthetic testbed behind the recommendations
(Section~\ref{sec:sel-testbed}). Section~\ref{sec:sel-benchmarks} states two standing benchmarks that any
revision of the pipelines must pass or explicitly fail; Section~\ref{sec:sel-causal} explains why the
synthetic worlds make the causal pipeline of Chapter~\ref{ch:causal} \emph{testable} rather than merely
arguable. Sections~\ref{sec:sel-uq}--\ref{sec:sel-limits} close with the uncertainty-quantification summary,
the production checklist, and the honest list of what this manual does not solve.

\section{Two decision tables}
\label{sec:sel-decisions}

There are two natural entry points. Either we start from the \emph{data situation} (what the vendor gave us,
which geography, how reliable the timestamps are) and ask what is estimable at all, or we start from the
\emph{question} and ask which pipeline answers it. Table~\ref{tab:sel-master} takes the first route;
Table~\ref{tab:sel-pipeline} the second. Both cite the chapter developing each recommendation; neither is a
substitute for reading it. The most common failure we observe in practice is running a question-driven
pipeline on a data situation that cannot support its estimand --- most often event-time excitation models on
bucketed exports (row three of Table~\ref{tab:sel-master}), or causal claims read off predictive fits (rows
five and six).

\begin{table}[htbp]
\centering
\footnotesize
\begin{tabular}{p{2.5cm}p{3.4cm}p{2.9cm}p{2.6cm}p{3.0cm}}
\toprule
Scenario & Recommended model & Estimand & UQ method & First-order risk \\
\midrule
EU fundamentals panel & feature-engineered XGBoost on point-in-time panel (Ch.~\ref{ch:xgb}) &
$\E[\text{success}\mid X]$; EBITDA-growth quantiles & split conformal / CQR & snapshot-field leakage \\
US marks, reliable day timestamps & marked Hawkes, ground $\times$ mark two-stage (Ch.~\ref{ch:hawkes}) &
excitation kernel; branching ratio & Laplace interval; parametric bootstrap & latent common factor
inflates excitation \\
US weekly or bucketed exports & covariate count GLM; MBPP interval-censored likelihood
(Ch.~\ref{ch:poisson},~\ref{ch:censored}) & covariate semi-elasticities; effective radius $\kappa$ &
quasi-Poisson sandwich; cluster bootstrap & aggregation manufactures apparent excitation \\
Short history or tiny sectors & pooled Poisson GLM with hierarchical shrinkage (Ch.~\ref{ch:poisson}) &
sector rate levels and trends & parametric bootstrap & overdispersion; a few large weeks dominate \\
Causal: financing effect on outcomes & backdoor adjustment; AIPW / DML (Ch.~\ref{ch:causal}) &
ATT of a round on survival / growth & cluster bootstrap CI; E-value & unobserved quality selects into
treatment \\
Causal: contagion between firms or sectors & event-time Hawkes with factor proxies and sensitivity
analysis (Ch.~\ref{ch:hawkes}) & excitation net of common shocks & profile CI; sensitivity band & radius
absorbs latent risk appetite \\
Macro-effect question & count GLM with publish-lag exposure design (Ch.~\ref{ch:us},~\ref{ch:poisson}) &
rate semi-elasticity to macro series & HAC / sandwich & publication-lag leakage; few macro cycles \\
\bottomrule
\end{tabular}
\caption{Master decision table: from data situation to method. Every row assumes the point-in-time
discipline of Chapter~\ref{ch:data} has already been enforced.}
\label{tab:sel-master}
\end{table}

\begin{table}[htbp]
\centering
\small
\begin{tabular}{p{4.4cm}p{9.8cm}}
\toprule
Question & Pipeline \\
\midrule
Success prediction & feature-engineered XGBoost \cite{chen2016} with conformal intervals
\cite{romano2019} on the point-in-time panel (Ch.~\ref{ch:xgb}) \\
Who-gets-funded ranking & two-stage sector count process plus funded-pool conditional logit
\cite{mcfadden1974} with an explicit newcomer option (Ch.~\ref{ch:hawkes}) \\
Contagion measurement & event-time Hawkes \cite{hawkes1971,bacry2015} with macro-factor proxies and
sensitivity analysis (Ch.~\ref{ch:hawkes}) \\
Macro sensitivity & covariate count GLM with exposure and publish-lag designs
(Ch.~\ref{ch:us},~\ref{ch:poisson}) \\
Censored or bucketed exports & MBPP interval-censored likelihood; report the effective radius $\kappa$,
never the raw kernel parameters $\theta$ (Ch.~\ref{ch:censored}) \\
\bottomrule
\end{tabular}
\caption{Question-driven decision table: from estimand to pipeline.}
\label{tab:sel-pipeline}
\end{table}

\section{Testbed evidence}
\label{sec:sel-testbed}

All quantitative claims below come from the synthetic testbed that ships with the code
(Appendix~\ref{app:code}), seed 0 throughout. Two worlds are simulated, both mirroring the schemas of
Chapter~\ref{ch:data}. \emph{Regime A} is well specified for our tools: a weekly two-stage DGP over 520
weeks, 12 sectors, and 8{,}000 firms (7{,}031 tracked), in which sector counts follow a count GLM with lagged
excitation and firms are selected by a risk-set softmax with a 26-week cooldown indicator (true coefficient
approximately $-2.7$); it produces 50{,}036 observed deals. \emph{Regime B} is deliberately misspecified: a
daily firm-level block Hawkes process with 30-day-half-life self-inhibition, 7-day-half-life sector
contagion, and a latent quality mark, producing 49{,}788 deals fitted with the same weekly tools. Both worlds
pass through an observation-noise layer calibrated to Chapter~\ref{ch:data}: 28\% undisclosed deal sizes, a
further 15\% only estimated (18\% relative error), 55\% missing valuations, 3\% sector mislabels, 1.2\%
duplicates, date jitter of up to 6 days, and 12\% of firms untracked. Regime A shows what the estimators
recover when their assumptions hold but the data are dirty; regime B shows what the same numbers mean when
the world is event-time and the export is not.

\subsection{Stage 1: sector-level count model recovery}
\label{sec:sel-stage1}

\begin{table}[htbp]
\centering
\small
\begin{tabular}{lcc}
\toprule
Quantity & Regime A & Regime B \\
\midrule
Covariate $\beta$ correlation with truth & 0.675 & 0.37 \\
Intercept bias & $-0.09$ & --- \\
Excitation--entry correlation & 0.27 & --- \\
Fitted effective spectral radius & 0.355 & 0.57 \\
True effective spectral radius & 0.172 & (not weekly-defined) \\
Held-out Pearson dispersion & --- & 2.72 \\
\bottomrule
\end{tabular}
\caption{Stage-1 recovery (seed 0). Regime A is well specified; regime B is a daily block Hawkes fitted
with weekly tools.}
\label{tab:sel-stage1}
\end{table}

Table~\ref{tab:sel-stage1} carries three lessons. First, observation noise alone attenuates: in regime A,
where the fitted class contains the truth, the covariate coefficients correlate 0.675 with the true $\beta$
--- roughly 25\% attenuation, consistent with errors-in-variables from estimated sizes and jittered dates.
Signs and rankings survive; magnitudes are lower bounds. Second, the fitted effective radius of 0.355
against a true 0.172 is \emph{not} noise: it is the latent serially correlated risk-appetite factor being
absorbed into the excitation term, exactly the confounding mechanism of Section~\ref{sec:sel-causal}. Third,
regime B shows what misspecification does: coefficient correlation collapses to 0.37, the fitted radius of
0.57 is mostly a daily-to-weekly aggregation artifact, and the held-out Pearson dispersion of 2.72 flags the
failure --- hence the dispersion check on the production checklist. What generalizes: the direction of every
bias (attenuation, radius inflation, dispersion under aggregation). What is synthetic-specific: the exact
factors, which depend on the noise calibration above.

\subsection{Stage 2: funded-pool conditional logit}
\label{sec:sel-choice}

\begin{table}[htbp]
\centering
\small
\begin{tabular}{lcc}
\toprule
Metric & Regime A & Regime B \\
\midrule
Held-out NLL per event & 3.346 & 3.796 \\
Random-choice NLL per event & 4.011 & 4.172 \\
Top-5 accuracy & 0.446 & 0.316 \\
Log-gap coefficient & $+1.146$ & $+0.437$ \\
Newcomer ASC & 6.856 & 4.691 \\
Newcomer share of events & $\approx 0.27$ & $\approx 0.27$ \\
Dropped immediate-repeat events & 35 & 2 \\
\bottomrule
\end{tabular}
\caption{Choice-model performance: 64-candidate sampled softmax on truthful risk sets, held-out weeks,
seed 0.}
\label{tab:sel-choice}
\end{table}

The model in Table~\ref{tab:sel-choice} is the funded-pool conditional logit with a 64-candidate sampled
softmax (consistency argument in Appendix~\ref{app:proofs}) on truthful risk sets. In regime A the recovered
feature weights are: raised-to-date 0.336 (truth 0.30), employees 0.253 (truth 0.30), age 0.054 (truth
$-0.15$), stage 0.132 (truth 0.15). Three of four are recovered to within attenuation; the age sign is
missed because the 26-week cooldown indicator and the fitted log-gap term ($+1.146$) absorb the age-adjacent
recency structure --- the mechanism is captured but attributed to timing features rather than to age. The
large newcomer alternative-specific constant is not a pathology: it prices the size of the unobserved
entrant pool, and by the integrated-entrant identity of Appendix~\ref{app:proofs} it is interpretable only
jointly with the assumed pool size. In regime B every number degrades but the ordering of models is
preserved, which is the property we actually rely on when ranking candidates.

\subsection{Risk-set contamination}
\label{sec:sel-riskset}

\begin{table}[htbp]
\centering
\small
\begin{tabular}{lcl}
\toprule
Risk-set construction $\RS_t$ & Held-out NLL & Weight behaviour \\
\midrule
Naive observed-active mask & 3.886 & inverted: recency $+1.94$ vs.\ true 0 \\
Oracle membership mask & 3.558 & all weights recovered \\
\bottomrule
\end{tabular}
\caption{Risk-set contamination, regime A. The only difference between rows is who is deemed at risk at
each event time.}
\label{tab:sel-riskset}
\end{table}

\begin{warning}
\label{warn:sel-riskset}
Risk-set validity is first-order. In Table~\ref{tab:sel-riskset} the naive mask --- treating a firm as at
risk only when the vendor happens to show it active --- does not merely cost likelihood; it \emph{inverts}
the recency weight to $+1.94$ against a true value of 0, because firms enter the observed mask by
transacting. No feature engineering repairs an estimate whose conditioning set is wrong. Audit $\RS_t$
before touching the utility specification.
\end{warning}

\subsection{A stability episode, and the operator surrogate}
\label{sec:sel-stability}

\begin{table}[htbp]
\centering
\small
\begin{tabular}{lccl}
\toprule
Feedback specification & Eff.\ radius & Outcome & Consequence \\
\midrule
Log-link, raw-count feedback & 0.35 & metastable & 50k-event world tipped to 169k events \\
Log-link, raw-count feedback & 0.15 & stable & simulation matches fitted moments \\
\bottomrule
\end{tabular}
\caption{Stability episode. Linear-Hawkes intuition (explosion only at radius $\ge 1$) does not transfer
to log-link feedback on raw counts.}
\label{tab:sel-stability}
\end{table}

Table~\ref{tab:sel-stability} records the incident behind checklist item 4: a log-link model feeding back
raw counts at effective radius 0.35 --- comfortably subcritical by linear intuition --- is metastable, and one
excursion tipped a 50{,}000-event world to 169{,}000 events. Stability theory for nonlinear self-excitation
\cite{bremaud1996,fokianos2009} makes clear that the linear threshold is not the relevant boundary once the
link is convex in the feedback; bounded transforms such as $\log(1+n)$ restore contractivity. Separately,
the PINO operator surrogate for the MBPP Volterra equation (derivation in Appendix~\ref{app:proofs}) reaches
sub-2\% relative $L^2$ error across 3--8 sector configurations under the full training protocol, at roughly
3\,ms per solve; a 600-epoch demonstration run reaches only 8.2\%, so the cheap configuration is a demo, not
a substitute.

\section{Standing benchmarks}
\label{sec:sel-benchmarks}

We maintain two permanent falsifiability devices. They are not experiments run once; they are challenges
every revision of the pipelines must confront, with targets frozen from Section~\ref{sec:sel-testbed}.

\textbf{Benchmark 1 (XGBoost ranking baseline).} Fit a gradient-boosted ranker \cite{chen2016} on
\emph{identical} choice sets --- same 64 candidates, same truthful risk sets, same temporal splits --- and
compare top-5 accuracy against the conditional logit's 0.446 (regime A) and 0.316 (regime B). Decision
rule: if the boosted ranker wins materially, the linear-in-features utility is missing features or
nonlinearity, and the remedy is feature work, not model exotica; if it does not, we keep the logit for its
interpretable weights and the ASC decomposition. Either outcome is informative; a run that skips the
identical-choice-set discipline is informative about nothing.

\textbf{Benchmark 2 (non-homogeneous Poisson null).} For any claimed excitation effect, fit the nested
ladder: covariates-only inhomogeneous Poisson (simulated by thinning \cite{lewis1979}), then covariates plus
excitation, then the full event-time timing model, and decompose the held-out improvement across the rungs.
If the excitation rung contributes approximately nothing at this data scale, drop it: an excitation term
that does not pay its way in held-out likelihood is a liability in simulation
(Table~\ref{tab:sel-stability}) and an attractive nuisance in interpretation (Table~\ref{tab:sel-stage1}).
The residual-analysis toolkit of \cite{ogata1988} supplies the diagnostics for each rung.

\section{Causal validation on synthetic ground truth}
\label{sec:sel-causal}

For a manager deciding whether to trust the causal layer, the testbed changes the conversation: because the
DGP is fully known, the pipeline of Chapter~\ref{ch:causal} is \emph{testable}, not merely defensible. The
protocol: (i) generate a binary treatment (say, receiving a growth round) with known selection on
observables \emph{and} on the latent quality mark; (ii) estimate the effect four ways --- naive difference in
means, backdoor-adjusted regression \cite{pearl1995}, AIPW, and double/debiased ML
\cite{imbens2015,chernozhukov2018} --- and compare each to the known truth; (iii) attach cluster bootstrap
confidence intervals \cite{efron1994}; (iv) compute E-values \cite{vanderweele2017} and check them against
the known strength of the latent confounding. The code asserts the expected qualitative outcomes as
hypotheses: the naive estimate is biased in the direction of selection; backdoor adjustment removes the
observable component; AIPW matches backdoor with better robustness to nuisance misspecification; and the
residual gap to truth is attributable to latent quality, with magnitude consistent with the computed
E-value. Where an assertion fails, the pipeline --- not the world --- is wrong.

The count-process analogue is already in evidence: the fitted effective radius of 0.355 against a true 0.172
(Table~\ref{tab:sel-stage1}) is the point-process version of a confounded treatment effect. ``Past events
excite future events'' is a causal claim in which past events are the treatment; the latent serially
correlated risk-appetite factor is the unmeasured common cause, and its omission violates exchangeability in
precisely the sense of \cite{hernan2020}. The excitation-vs-latent-factor experiment therefore doubles as
the standing reminder that no in-sample fit statistic distinguishes contagion from common shocks; only
design (factor proxies, sensitivity bands, and, on the simulator, $\Do$-style interventions) does.

\section{Uncertainty quantification: summary}
\label{sec:sel-uq}

\begin{table}[htbp]
\centering
\footnotesize
\begin{tabular}{p{2.9cm}p{2.9cm}p{3.6cm}p{4.4cm}}
\toprule
Method & Technique & Guarantee & Key assumptions \\
\midrule
Count GLM (Ch.~\ref{ch:poisson}) & quasi-Poisson sandwich &
asymptotic SE validity for mean parameters & correct mean function; dispersion as nuisance \\
Two-stage / choice (Ch.~\ref{ch:hawkes}) & cluster bootstrap by week &
asymptotic CI coverage & weak dependence across clusters; clustering at the right level \\
XGBoost prediction (Ch.~\ref{ch:xgb}) & split conformal / CQR \cite{romano2019} &
finite-sample marginal coverage & exchangeability; weakened under temporal splits \\
Hawkes MLE (Ch.~\ref{ch:hawkes}) & Laplace / posterior interval &
local asymptotic normality & correct specification; stationarity \\
MBPP (Ch.~\ref{ch:censored}) & parametric bootstrap &
coverage under the fitted DGP & model correctness; report $\kappa$, not $\theta$ \\
\bottomrule
\end{tabular}
\caption{UQ summary: what each interval means and what it silently assumes.}
\label{tab:sel-uq}
\end{table}

Table~\ref{tab:sel-uq} is deliberately blunt about guarantees. Only the conformal row carries a
finite-sample statement, and even that one is marginal (not conditional) and degrades when temporal splits
break exchangeability. Every other interval in this manual is asymptotic and specification-dependent; the
dispersion of 2.72 in Table~\ref{tab:sel-stage1} means naive Poisson standard errors in regime B are
understated by about $\sqrt{2.72} \approx 1.65$ before correction.

\section{Production checklist}
\label{sec:sel-checklist}

\begin{enumerate}
\item \textbf{Point-in-time audit.} Reconstruct every feature from data knowable at the prediction date;
join macro series on \texttt{publish\_date}, never \texttt{ref\_date} (Ch.~\ref{ch:data},~\ref{ch:us}).
\item \textbf{Deduplication.} Apply the 3-day-window convention before any counting; the testbed carries
1.2\% duplicates and date jitter up to 6 days, and both alias into spurious short-lag excitation.
\item \textbf{Risk-set truthfulness audit.} Validate the membership mask $\RS_t$ against a labelled
subsample of known entrants and exits; contamination inverts weights (Table~\ref{tab:sel-riskset}).
\item \textbf{Effective-radius guard.} Refuse to simulate any fitted model with effective radius $\ge 0.95$;
for log-link raw-count feedback refuse well below that (metastability observed at 0.35,
Table~\ref{tab:sel-stability}), or use bounded transforms.
\item \textbf{Dispersion correction.} Report the held-out Pearson dispersion and scale standard errors
accordingly; treat dispersion far above 1 as a specification alarm, not a nuisance.
\item \textbf{Universe coverage and newcomer calibration.} Track the untracked-firm share (12\% in the
testbed) and recalibrate the newcomer ASC whenever the pool definition changes; the ASC is meaningless
detached from the pool (Appendix~\ref{app:proofs}).
\item \textbf{Temporal splits only.} No random splits anywhere in the stack --- not for model selection, not
for conformal calibration, not for the benchmarks of Section~\ref{sec:sel-benchmarks}.
\item \textbf{Seeds and reproducibility.} Every reported number carries a seed (this chapter: seed 0) and a
regeneration command; Appendix~\ref{app:code} maps each table here to its entry point.
\end{enumerate}

\section{Limitations and roadmap}
\label{sec:sel-limits}

Four limitations stand. First, entity resolution is single-vendor: without a second source,
\texttt{company\_id} merges and splits cannot be audited, and every firm-level count inherits that
fragility. Second, latent-factor confounding is mitigated, not solved: the 0.355-versus-0.172 radius gap
persists under our best factor proxies, and we know of no observational design that closes it; sensitivity
reporting is the honest ceiling. Third, cross-border generalization is untested: the two synthetic worlds
mirror the EU-fundamentals and US-marks regimes respectively, and nothing here licenses transporting a model
fitted in one to the other. Fourth, the outcome space is too coarse: ``next event'' conflates a follow-on
round with an acquisition and with quiet failure. The next modelling step is competing risks ---
cause-specific hazards for M\&A versus next round in the spirit of \cite{cox1972}, marked into the
point-process machinery of \cite{daley2003} --- for which the survival scaffolding in the code base
(Appendix~\ref{app:code}) is already laid.
